% ABSTRACT (English)
% Note: Should not exceed 300 words.
% Should cover: Objective, Methodology, Results, Applications
% Use Past tense for objectives and methodology
% Use Present tense for results and applications
%
This project aims to design and develop a mobile application named ``Pacegasus'' to promote sustainable running behavior. The system integrates Gamification principles with an Adaptive Training Engine and a social interaction framework, grounded in Self-Determination Theory (SDT) and Behavior Change Techniques (BCTs).

The application is developed using React Native for cross-platform Android deployment, leveraging built-in smartphone sensors---GPS and Accelerometer---to track distance, average pace, and step count without requiring external wearable devices. The backend is powered by Node.js and PostgreSQL for data processing and storage.

Key features include an Adaptive Training Engine that analyzes daily wellness check-in data, Session-RPE scores, and the Acute-to-Chronic Workload Ratio (ACWR) to personalize daily and weekly missions. Gamification elements such as points, badges, leaderboards, and a Story Mode are incorporated alongside social features enabling users to add friends, join group activities, and share running results.

The expected outcome is a functional prototype capable of improving exercise adherence, reducing injury risk, and fostering long-term healthy behaviors---contributing to the mitigation of Non-Communicable Diseases (NCDs) at the population level.
